\section{Conclusion}
\label{sec:conclusion}

Making patch-level text--image alignment an explicit training objective turns a
globally-aligned vision--language model into a text-prompted organ localizer,
recovering most of the localization signal for morphologically distinctive
structures without any per-pixel label at inference. The same frozen model
transfers to an unseen dataset and acquisition protocol to answer a downstream
clinical question, ranking the organ a vertebral metastasis is nearest to
invading well above chance across 180 lesions.

Two findings from that pipeline generalize past it. The first is that prompt
elaboration is not what makes text supervision work, because the class name
carries almost all of the signal and multi-sentence anatomical description adds
nothing measurable over a single generic sentence. The second is that when a
text-prompted localizer fails, it fails confidently, segmenting the wrong
structure cleanly rather than signaling uncertainty, which is why prompt
validation belongs in any deployed system as a first-class component.

Against those findings, our method remains bounded by the $14\times14$ patch
grid, by a two-class failure that we diagnose but do not repair, and by
single-cohort evaluation of the downstream index. The next experiment is
therefore specified by the diagnosis rather than chosen freely: we should scope
hard negatives to the classes they were introduced to correct, which is the one
change the failed global-negative result actually implicates.
